\documentclass[11pt]{article}
\usepackage{amsmath,amsfonts,mathabx}
\usepackage{graphicx}
\usepackage{hyperref}

% PREAMBLE
\newcommand{\RR}{\ensuremath{\mathbb{R}}}
\newcommand{\RRn}[1]{\mathbb{R}^{#1}}
\DeclareMathOperator{\function}{function}

% Code listing
\usepackage[listings]{tcolorbox}
\usepackage[normalem]{ulem}
\usepackage{xcolor}


\begin{document}


\begin{center}
{\textbf{\LARGE Chapitre 3: Pr\'ediction, taux d'erreur, s\'election de mod\`eles et r\'egularisation}}
\end{center}

\tableofcontents
\newpage

\section{Taux d'erreur}

\noindent Pour faciliter la compr\'ehension des notions pr\'esent\'ees, on limite la mati\`ere pour le moment \`a la r\'egression lin\'eaire. Les g\'en\'eralisations \`a la r\'egression logistique seront abord\'ees plus tard.

\noindent Consid\'erons le jeu de donn\'ees \textsf{Auto} de la library \texttt{ISLR} du logiciel \textrm{R}. Ce jeu de donn\'ees contient 392 observations. La variable r\'eponse est \texttt{mpg}, la distance moyenne parcourue (en miles) par la voiture avec un gallon de carburant. On s'int\'eresse \`a une seule variable explicative, soit la puissance de la voiture \textrm{horsepower}.

\noindent On veut choisir un parmi les mod\`eles suivants: 
\begin{eqnarray*}
\mbox{Mod\`ele 1: } Y_i &=& \beta_0 + \beta_1 X_i + \epsilon_i  \\
\mbox{Mod\`ele 2: } Y_i &=& \beta_0 + \beta_1 X_i + \beta_2 \log(X_i) + \epsilon_i \\
\mbox{Mod\`ele 3: } Y_i &=& \beta_0 + \beta_1 X_i + \beta_2 X_i^2 + \epsilon_i \\
\mbox{Mod\`ele 4: } Y_i &=& \beta_0 + \beta_1 X_i + \beta_2 X_i^2 + \beta_3 X_i^3 + \epsilon_i \\
\end{eqnarray*}
\noindent Pour tous ces mod\`eles, l'indice $i$ varie de $1$ \`a $n=392$ et on suppose que $\epsilon_i \sim {\cal N}(0,\sigma^2)$.

\noindent Pour chaque mod\`ele, on estime les coefficients $\beta_j$ par la m\'ethode des moindres carr\'es et on calcule les valeurs pr\'edites $\hat Y_i, i=1,\cdots,n.$ Par exemple, pour le mod\`ele 1, on a $\hat Y_i = \hat\beta_0+\hat\beta_1 X_i$. Par la suite, on calcule la somme des carr\'es des r\'esidus $$SSE=\sum\limits_{i=1}^n (Y_i-\hat Y_i)^2.$$

\noindent Pour choisir entre les 4 mod\`eles list\'es plus haut, on peut calculer des crit\`eres comme l'$AIC$ ou le $BIC$ ou effectuer des tests d'hypoth\`eses pour les mod\`eles embo\^it\'es. Toutes ces proc\'edures reviennent \`a comparer des $SSE$s des diff\'erents mod\`eles. Ces m\'ethodes cherchent donc \`a identifier le mod\`ele avec le meilleur ajustement aux donn\'ees.

\noindent  Supposer maintenant qu'on a une nouvelle valeur de la variable explicative $X_0$ et on veut s\'electionner le mod\`ele qui pr\'edit la variable r\'eponse $Y_0$ avec la meilleure pr\'ecision. Autrement dit, on veut identifier le mod\`ele qui minimise $E[(Y_0-\hat Y_0)^2]$. 

\noindent Les deux crit\`eres (ad\'equation et pr\'ecision de la pr\'ediction) sont diff\'erents et peuvent m\^eme \^etre contradictoires. Si on a $n$ points $(X_i,Y_i)$, un mod\`ele polynomial d'ordre $n-1$ aura une ad\'eqaution parfaite ($SSE=0$). Cependant, ce mod\`ele ne sera pas utile si on veut pr\'edire des valeurs futures.

\noindent L'objectif de ce chapitre est pr\'esenter des m\'ethodes pour s\'electionner des mod\`eles en se basant sur la pr\'ecision de pr\'ediction. 

\noindent \`A cet effet, il est important de distinguer deux types d'erreur: taux d'erreur entrainement et taux d'erreur test.

\noindent Le taux d'erreur entrainement peut \^etre \'evalu\'e par $SSE$. Il mesure l'ad\'equation des mod\`eles aux donn\'ees. On l'utilise pour comparer l'ad\'equation de plusieurs mod\`eles \`a un jeu de donn\'ees.

\noindent Le taux d'erreur test \'evalue la pr\'ecision d'une observation future. Dans la prochaine section, on pr\'esente quelques m\'ethodes pour l'estimer.

\section{Validation crois\'ee}

\noindent L'id\'ee de la validation crois\'ee est de partager le jeu de donn\'ees en deux parties. La premi\`ere partie est appel\'ee jeu de donn\'ees entrainement. Elle sert \`a estimer les param\`etres du mod\`ele. La deuxi\`eme partie est le jeu de donn\'ees test. Elle sert \`a estimer le taux d'erreur test.

\noindent Revenons au jeu de donn\'ees \textsf{Auto} de la library \texttt{ISLR} du logiciel \textrm{R}. Il est compos\'e de $n=392$ observations. On le partage en deux jeux de donn\'ees. Le premier est compos\'e des $n_1=275$ premi\`eres observations. Il sert \`a estimer les param\`etres $\beta$. Le deuxi\`eme est compos\'es des $n_2=n-n_1=392-275=117$ derni\`eres observations. Il sert \`a estimer le taux d'erreur test par $$\frac{1}{n_2}\sum\limits_{i=n_1+1}^{\sout{n_2}\textcolor{red}{n}} (Y_i-\hat Y_i)^2=\frac{1}{117} \sum\limits_{i=276}^{392} (Y_i-\hat Y_i)^2.$$ 
\noindent Si on consid\`ere le mod\`ele 1, alors $\hat Y_i=\hat\beta_0+\hat\beta_1 X_{i}$ est la valeur pr\'edite de $Y_i$ avec $\hat\beta_0$ et $\hat\beta_1$ calcul\'e \`a partir de $\{(Y_i,X_{i}), i=1,\cdots,275\}$.

\begin{verbatim}
> n <- 392
> n1 <- 275
> entrainement <- 1:n1
> modele.1.entrainement <- lm(mpg ~ horsepower,
                              data=Auto,subset=entrainement)
> modele.2.entrainement <- lm(mpg ~ horsepower+log(horsepower),
                              data=Auto,subset=entrainement)
> modele.3.entrainement <- lm(mpg ~ horsepower+I(horsepower^2),
                              data=Auto,subset=entrainement)
> modele.4.entrainement <- lm(mpg ~ horsepower+I(horsepower^2)+I(horsepower^3),
                              data=Auto,subset=entrainement)
> 
> test <- (n1+1):n
> Y.chapeau.1 <- predict(modele.1.entrainement,newdata=Auto[test,])
> Y.chapeau.2 <- predict(modele.2.entrainement,newdata=Auto[test,])
> Y.chapeau.3 <- predict(modele.3.entrainement,newdata=Auto[test,])
> Y.chapeau.4 <- predict(modele.4.entrainement,newdata=Auto[test,])
> 
> Y <- Auto$mpg
> 
> err.test.modele.1 <- mean((Y[test]-Y.chapeau.1)^2) ; err.test.modele.1
[1] 64.50982
> err.test.modele.2 <- mean((Y[test]-Y.chapeau.2)^2) ; err.test.modele.2
[1] 53.69829
> err.test.modele.3 <- mean((Y[test]-Y.chapeau.3)^2) ; err.test.modele.3
[1] 52.9343
> err.test.modele.4 <- mean((Y[test]-Y.chapeau.4)^2) ; err.test.modele.4
[1] 52.7879
\end{verbatim}

\noindent C'est le mod\`ele 4 qui semble avoir la meilleure pr\'ediction. Cependant, la partition du \sout{jeux}\textcolor{red}{jeu} de donn\'ees original en deux est arbitraire. La fa\c{c}on la plus facile de r\'eduire ce al\'ea d\^u \`a la r\'epartition du \sout{jeux}\textcolor{red}{jeu} de donn\'ees en deux est de recommencer cette proc\'edure un grand nombre de fois et de prendre des moyennes. 

\begin{verbatim}
> if (!require("ISLR")) install.packages("ISLR")
> library("ISLR")#### Installer la librairie ISLR
> 
> data(Auto)
> 
> n <- 392
> n1 <- 275
> 
> B <- 1000
> err.test.modele.1 <- NULL
> err.test.modele.2 <- NULL
> err.test.modele.3 <- NULL
> err.test.modele.4 <- NULL
> 
> for(b in 1:B)
+ {
+ entrainement <- sample(n,n1)
+ 
+ modele.1.entrainement <- lm(mpg ~ horsepower,data=Auto,subset=entrainement)
+ modele.2.entrainement <- lm(mpg ~ horsepower+log(horsepower),
+                             data=Auto,subset=entrainement)
+ modele.3.entrainement <- lm(mpg ~ horsepower+I(horsepower^2),
+                             data=Auto,subset=entrainement)
+ modele.4.entrainement <- lm(mpg ~ horsepower+I(horsepower^2)+I(horsepower^3),
+                             data=Auto,subset=entrainement)
+ 
+ Y.chapeau.1 <- predict(modele.1.entrainement,newdata=Auto[-entrainement,])
+ Y.chapeau.2 <- predict(modele.2.entrainement,newdata=Auto[-entrainement,])
+ Y.chapeau.3 <- predict(modele.3.entrainement,newdata=Auto[-entrainement,])
+ Y.chapeau.4 <- predict(modele.4.entrainement,newdata=Auto[-entrainement,])
+ 
+ Y <- Auto$mpg
+ 
+ err.test.modele.1 <- c(err.test.modele.1,mean((Y[-entrainement]-Y.chapeau.1)^2)) 
+ err.test.modele.2 <- c(err.test.modele.2,mean((Y[-entrainement]-Y.chapeau.2)^2)) 
+ err.test.modele.3 <- c(err.test.modele.3,mean((Y[-entrainement]-Y.chapeau.3)^2)) 
+ err.test.modele.4 <- c(err.test.modele.4,mean((Y[-entrainement]-Y.chapeau.4)^2)) 
+ }
> 
> 
> mean(err.test.modele.1)
[1] 24.27115
> mean(err.test.modele.2)
[1] 19.81173
> mean(err.test.modele.3)
[1] 19.30965
> mean(err.test.modele.4)
[1] 19.42957
\end{verbatim}

\noindent En fin de compte, le mod\`ele 3 semble \^etre le meilleur pour pr\'edire des observations futures. M\^eme si on r\'ep\`ete l'op\'eration plusieur\textcolor{red}{s} fois, on obtient le m\^eme r\'esultat.

\noindent Un deuxi\`eme d\'efaut des deux approches pr\'esent\'ees est la pr\'ecision des estimateurs $\hat \beta$. En effet, nous sommes en train d'estimer $\beta$ avec seulement $70\%$ des donn\'es. Et donc la variance de l'estimateur $\hat\beta$ est relativement grande. 

\noindent Pour rem\'edier \`a cela, pour $i=\textcolor{red}{1},\sout{,}\cdots,n$, on estime $\hat\beta$ avec toutes les observations sauf la $i^{\mbox{\`eme}}$. On calcule $\hat y_{i,(-i)}$. Le crit\`ere sera donc $$VC_{(n)} = \frac{\sum\limits_{i=1}^n (y_i-\hat y_{i,(-i)})^2}{n}.$$

\noindent On montre que $$VC_{(n)} = \frac{1}{n}  \sum\limits_{i=1}^n \left(\frac{y_i-\hat y_i}{1-h_i}\right)^2,$$

\noindent o\`u $h_i$ est le levier de l'observation $i$. Et par cons\'equent, on n'a pas besoin d'ajuster $n$ r\'egressions lin\'eaires. Le code informatique est plus simple et surtout plus rapide.
\begin{verbatim}
> if (!require("ISLR")) install.packages("ISLR")
> library("ISLR")#### Installer la librairie ISLR
> 
> data(Auto)
> 
> CV.n <- function(modele)
+ {
+ residus <- modele$residuals
+ h <- hatvalues(modele)
+ resultat <- mean((residus/(1-h))^2)
+ return(resultat)
+ }
> 
> modele.1 <- lm(mpg ~ horsepower,data=Auto)
> modele.2 <- lm(mpg ~ horsepower+log(horsepower),
+                             data=Auto)
> modele.3 <- lm(mpg ~ horsepower+I(horsepower^2),
+                             data=Auto)
> modele.4 <- lm(mpg ~ horsepower+I(horsepower^2)+I(horsepower^3),
+                             data=Auto)
> 
> CV.n(modele.1)
[1] 24.23151
> CV.n(modele.2)
[1] 19.70068
> CV.n(modele.3)
[1] 19.24821
> CV.n(modele.4)
[1] 19.33498
\end{verbatim}

\noindent Le mod\`ele 3 est encore une fois le meilleur pour pr\'edire des observations futures.

\noindent Le calcul de $VC_{(n)}$ est direct dans le cas d'une r\'egression lin\'eaire. Cependant, pour d'autres mod\`eles de r\'egression dans la famille des r\'e\textcolor{red}{g}ressions lin\'eaires g\'e\textcolor{red}{n}\'eralis\'e\textcolor{red}{e}s, ce calcul requi\`ere l'ajustement de $n$ mod\`eles de r\'egression, ce qui peut \^etre difficile d'un point de vue informatique. Pour cette raison, on peut consid\'erer une autre version de validation crois\'ee $VC_{(K)}$. Dans la pratique, on utilise $K=5$ ou $K=10$. On partage le jeu de donn\'ees en $K$ parties de tailles quasiment \'egales: $A_1,\cdots,A_K$. Pour $k=1,\cdots,K$, on estime les param\`etres $\beta$ \`a partir de toutes les donn\'ees sauf $A_k$ et on utilise $A_k$ pour estimer l'erreur avec
$$MSE_k = \frac{1}{n_k} \sum_{i \in A_k} \left\{y_i-\hat y_{i,(-k)}\right\}^2,$$
o\`u $n_k$ est la taille de $A_k$. On estime par la suite l'erreur test par :
$$VC_{(K)} = \frac{1}{K}  \sum\limits_{k=1}^K MSE_k.$$

\begin{verbatim}
> ###########################################################
> ## Calcul de CV.n et CV.k en utilisant la librairie boot ##
> ###########################################################
> 
> if (!require("boot")) install.packages("boot")
> library("boot")#### Installer la librairie boot
> 
> modele.glm.1 <- glm(mpg ~ horsepower,data=Auto)
> modele.glm.2 <- glm(mpg ~ horsepower+log(horsepower),
+                             data=Auto)
> modele.glm.3 <- glm(mpg ~ horsepower+I(horsepower^2),
+                             data=Auto)
> modele.glm.4 <- glm(mpg ~ horsepower+I(horsepower^2)+I(horsepower^3),
+                             data=Auto)
> 
> ## CV.n
> cv.glm(Auto, modele.glm.1)$delta[1]
[1] 24.23151
> cv.glm(Auto, modele.glm.2)$delta[1]
[1] 19.70068
> cv.glm(Auto, modele.glm.3)$delta[1]
[1] 19.24821
> cv.glm(Auto, modele.glm.4)$delta[1]
[1] 19.33498
> 
> ## CV.K (K=5)
> cv.glm(Auto, modele.glm.1, K=5)$delta[1]
[1] 24.26099
> cv.glm(Auto, modele.glm.2, K=5)$delta[1]
[1] 19.80128
> cv.glm(Auto, modele.glm.3, K=5)$delta[1]
[1] 19.24542
> cv.glm(Auto, modele.glm.4, K=5)$delta[1]
[1] 19.30565
> 
> ## CV.K (K=10)
> cv.glm(Auto, modele.glm.1, K=10)$delta[1]
[1] 24.13985
> cv.glm(Auto, modele.glm.2, K=10)$delta[1]
[1] 19.92659
> cv.glm(Auto, modele.glm.3, K=10)$delta[1]
[1] 19.23274
> cv.glm(Auto, modele.glm.4, K=10)$delta[1]
[1] 19.46304
\end{verbatim}

\section{R\'egression r\'egularis\'ee}

Consid\'erons un mod\`ele de r\'egression
$$Y_i = \beta_0 + \beta_1 X_{i1} + \cdots  + \beta_p X_{ip} + \epsilon_i; \, \, \, i=1,\cdots,n; \, \, \, \epsilon_i \sim {\cal N}(0,\sigma^2).$$

\noindent On estime $\beta=(\beta_0,\beta_1,\cdots,\beta_p)^\top$ en minimisant 
$$l(\beta_0,\beta_1,\cdots,\beta_p) = \sum_{i=1}^n \left\{Y_i - (\beta_0 + \beta_1 X_{i1} + \cdots  + \beta_p X_{ip})\right\}^2.$$

\noindent En \'ecriture matricielle, on a $l(\beta)=(Y-X\beta)^\top(Y-X\beta)$. L'estimateur obtenu v\'erifie $$\hat\beta \sim {\cal N}\left\{\beta,\sigma^2 (X^\top X)^{-1}\right\}.$$ 

\noindent On montre que $\hat\beta$ est sans biais et de variance minimale parmi les estimateurs sans biais. Cependant, rien de garantie que $\hat\beta$ poss\`ede la plus petite erreur quadratique moyenne. Pour $j=1,\cdots,p$, l'erreur quadratique moyenne de $\beta_j$ est $$E\left[(\beta_j-\hat\beta_j)^2\right]=\left[E(\hat\beta_j)-\beta_j\right]^2+Var(\hat\beta_j) = \mbox{biais}(\hat\beta_j)^2 + Var(\hat\beta_j).$$

\noindent La r\'egression r\'egularis\'ee impose une p\'enalit\'e \`a $l(\beta_0,\beta_1,\cdots,\beta_p)$. On estime alors $\beta$ en minimisant

$$l_\lambda(\beta_0,\beta_1,\cdots,\beta_p) = l(\beta_0,\beta_1,\cdots,\beta_p) + \lambda g(\beta_1,\cdots,\beta_p),$$

\noindent o\`u $g$ est une fonction de p\'enalit\'e connue et $\lambda$ est un param\`etre, appel\'e coefficient de p\'enalit\'e, qui permet de contr\^oler l'impact de la p\'enalit\'e impos\'ee \`a $\hat\beta$. L'objectif de cette p\'enalit\'e est de r\'eduire la variance des estimateurs obtenus. Cependant, cette p\'enalit\'e introduit un biais. Avec un choix judicieux de la fonction $g$ et du param\`etre $\lambda$, on diminue l'erreur quadratique moyenne des estimateur et on obtient des mod\`eles avec pr\'edictions plus pr\'ecises. Dans les sous-sections suivantes, on pr\'esente les r\'egressions r\'egularis\'ees les plus populaires.



\subsection{R\'egression Ridge}


\noindent La r\'egression ridge estime $\beta$ en minimisant
$$l^{(R)}_\lambda(\beta_0,\beta_1,\cdots,\beta_p) = \sum_{i=1}^n \left\{Y_i - (\beta_0 + \beta_1 X_{i1} + \cdots  + \beta_p X_{ip})\right\}^2 + \lambda \sum_{j=1}^p \beta_j^2.$$

\noindent En \'ecriture matricielle, on a $l^{(R)}_\lambda(\beta)=(Y-X\beta)^\top(Y-X\beta) + \lambda \beta^\top \tilde I_p \beta$. Ici, $\tilde I_p$ est une matrice diagonale de taille $p+1$. Le premier \'el\'ement de la diagonale de $\tilde I_p$ est \'egal \`a 0. Tous les autres \'el\'ements de la diagonale de $\tilde I_p$ sont \'egaux \`a 1. On obtient alors:
$$\hat\beta^{(R)}_{\lambda} = (X^\top X + \lambda \tilde I_p)^{-1} X^\top Y$$
\noindent et on en d\'eduit que
\begin{eqnarray*}
E[\hat\beta^{(R)}_{\lambda}] &=& (X^\top X + \lambda \tilde I_p)^{-1} X^\top X \beta  \neq \beta\\
Var[\hat\beta^{(R)}_{\lambda}] &=&  \sigma^2 (X^\top X + \lambda \tilde I_p)^{-1} X^\top X (X^\top X + \lambda \tilde I_p)^{-1}
\end{eqnarray*}
\noindent L'estimateur $\hat\beta^{(R)}_{\lambda}$ est donc biais\'e, sauf si $\lambda=0$. 

\noindent Consid\'erons maintenant le jeu de donn\'ees \textsf{Hitters} de la library \texttt{ISLR} du logiciel \textrm{R}. La variable r\'eponse est le salaire d'un joueur professionnel de base-ball au d\'ebut de la saison 1987 en milliers de dollars. Les variables explicatives portent essentiellement sur les performances et les caract\'eristiques des joueurs. On consid\`ere un mod\`ele lin\'eaire avec toutes les variables explicatives, sans interactions, et on estime les param\`etres $\beta_j, j=1,\cdots,19$  en utlisant une r\'egression ridge avec diff\'erentes valeurs de $\lambda$.

\begin{verbatim}
> if (!require("ISLR")) install.packages("ISLR")
> library("ISLR")#### Installer la librairie ISLR
> 
> data(Auto)
> 
> if (!require("glmnet")) install.packages("glmnet")
> library("glmnet") #### Installer la librairie glmnet
> 
> #############################
> ## Preparation des donnees ##
> #############################
> 
> data(Hitters)
> Hitters <- Hitters[!is.na(Hitters$Salary),]
> x <- model.matrix(Salary ~ ., Hitters)[, -1]
> y <- Hitters$Salary
> 
> #########################################################
> ## Regression Ridge avec differentes valeurs de lambda ##
> #########################################################
> 
> reg.ridge <- glmnet(x, y, alpha = 0)
> beta.chapeau.ridge <- t(coef(reg.ridge)[-1,])
> lambda.ridge <- reg.ridge$lambda
> log.lambda.ridge <- log(lambda.ridge)
> matplot(log.lambda.ridge,beta.chapeau.ridge,type="l",
+         xlab="Lambda (a l'echelle logarithmique)",ylab="beta chapeau (ridge)")
> abline(h=0)
\end{verbatim}

\noindent Et on obtient

\begin{center}
\includegraphics[scale=0.3]{Graphique1.pdf}
\end{center}

\noindent Lorsque $\lambda=0$, on a $\hat\beta^{(R)}_{\lambda}=\hat\beta$ et lorsque $\lambda \rightarrow \infty$, les estimateurs de $\beta_j, j=1,\cdots,p$ tendent tous vers 0 et l'estimateur de $\hat\beta^{(R)}_{0,\lambda}$ tend vers la moyenne $\bar{Y}$. Pour illustrer l'effet de la r\'egression ridge, on fixe $\lambda=e^{6}$ et on compare $\sum\limits_{i=1}^p \hat\beta_j^2$ \`a $\sum\limits_{i=1}^p [\hat\beta^{(R)}_{j,\lambda}]^2.$

\begin{verbatim}
> lambda.ridge <- exp(6)
> reg.ridge.2 <- glmnet(x, y, alpha = 0, lambda=lambda.ridge)
> reg.ridge.3 <- glmnet(x, y, alpha = 0, lambda=0)
> 
> sum((coef(reg.ridge.2)[-1])**2)
[1] 5720.707
> sum((coef(reg.ridge.3)[-1])**2)
[1] 18468.89
\end{verbatim}

\noindent On trouve $\sum\limits_{i=1}^p \hat\beta_j^2=18468.89$ et $\sum\limits_{i=1}^p [\hat\beta^{(R)}_{j,\lambda}]^2=5720.707$. On parle\sout{nt} d'un ph\'enom\`ene de r\'etr\'ecissement.

\noindent G\'en\'eralement, la r\'egression ridge diminue consid\'erablement le taux d'erreur test. Par cons\'equent, elle \sout{cro\^it}\textcolor{red}{am\'eliore} la pr\'ecision des pr\'edictions. Estimons le taux d'erreur test avec $\lambda=e^6$ et comparons ce taux d'erreur avec celui de la r\'egression lin\'eaire standard, qui correspond \`a $\lambda=0$.

\begin{verbatim}
> x.entrainement <- x[1:200,]
> y.entrainement <- y[1:200]
> reg.ridge.2.entrainement <- glmnet(x.entrainement,y.entrainement,
+                                    alpha=0,lambda=lambda.ridge)
> 
> x.test <- cbind(rep(1),x[201:263,])
> y.test <- y[201:263]
> y2.predites <- x.test%*%as.vector(coef(reg.ridge.2.entrainement))
> err.test.2 <- mean((y.test-y2.predites)**2); err.test.2
[1] 91274.42
> 
> reg.ridge.3.entrainement <- glmnet(x.entrainement,y.entrainement,
+                                    alpha=0,lambda=0)
> y3.predites <- x.test%*%as.vector(coef(reg.ridge.3.entrainement))
> err.test.3 <- mean((y.test-y3.predites)**2); err.test.3
[1] 124716.2
\end{verbatim} 

\noindent La diminution du taux d'erreur test est consid\'erable (pr\`es de $27\%$).

\subsection{R\'egression Lasso}

La r\'egression Lasso estime $\beta$ en minimisant

$$l^{(L)}_\lambda(\beta_0,\beta_1,\cdots,\beta_p) = \sum_{i=1}^n \left\{Y_i - (\beta_0 + \beta_1 X_{i1} + \cdots  + \beta_p X_{ip})\right\}^2 + \lambda \sum_{j=1}^p |\beta_j|.$$

\noindent Il n'existe pas d'expression explicite pour l'estimateur obtenu $\hat\beta^{(L)}_{\lambda}$. Cependant, il existe de nos jours plusieurs algorithms d'optimisation pour minimiser $l^{(L)}_\lambda(\beta_0,\beta_1,\cdots,\beta_p)$.

\noindent Une carat\'eristique int\'eressante de la r\'egression Lasso est que pour les grandes valeurs de $\lambda$, certains terme $\hat\beta_j$ deviennent identiquement \'egals \`a 0. Cet aspect permet une m\'ethode tr\`es efficace de s\'election de mod\`ele. 

\noindent Encore une fois, lorsque $\lambda=0$, on a $\hat\beta^{(L)}_{\lambda}=\hat\beta$. Lorsque $\lambda$ croit, le biais de $\hat\beta^{(L)}_{\lambda}$ \sout{diminue}\textcolor{red}{cro\^it} alors que sa variance \sout{cro\^it}\textcolor{red}{diminue}.

\begin{verbatim}
> reg.lasso <- glmnet(x, y, alpha = 1)
> beta.chapeau.lasso <- t(coef(reg.lasso)[-1,])
> lambda.lasso <- reg.lasso$lambda
> log.lambda.lasso <- log(lambda.lasso)
> matplot(log.lambda.lasso,beta.chapeau.lasso,type="l",
+         xlab="Lambda (a l'echelle logarithmique)",ylab="beta chapeau (lasso)")
> abline(h=0)
\end{verbatim}

\noindent Et on obtient

\begin{center}
\includegraphics[scale=0.3]{Graphique2.pdf}
\end{center}

\noindent On se fixe $\lambda=e^2$ et on calcule $\hat\beta^{(L)}_{\lambda}$. 

\begin{verbatim}
> ## On se fixe un lambda : log(lambda)=2
> 
> lambda.lasso <- exp(2)
> 
> reg.lasso.2 <- glmnet(x, y, alpha = 1, lambda=lambda.lasso)
> coef(reg.lasso.2)
20 x 1 sparse Matrix of class "dgCMatrix"
                       s0
(Intercept)   16.56691118
AtBat         -0.28054898
Hits           2.83421809
HmRun          .         
Runs           .         
RBI            .         
Walks          2.43602385
Years         -2.17005788
CAtBat         .         
CHits          .         
CHmRun         0.07794600
CRuns          0.25027060
CRBI           0.41108058
CWalks        -0.01322173
LeagueN       24.41162017
DivisionW   -117.77228136
PutOuts        0.24305524
Assists        .         
Errors        -0.89001203
NewLeagueN     .         
\end{verbatim}

\noindent Et on trouve que plusieurs coefficients $\hat\beta^{(L)}_{j,\lambda}$ sont identiquement \'egaux \`a z\'ero. Ceci nous pro\sout{d}\textcolor{red}{du}re une m\'ethode intuitive et naturelle de s\'election de mod\`eles.

\subsection{S\'election du param\`etre $\lambda$}

Il existe quelques m\'ethodes pour s\'electionner le param\`etre $\lambda$. La m\'ethode la plus intuitive est celle bas\'ee sur le taux d'erreur test. Ce taux d'erreur est estim\'e par validation crois\'ee pour plusieurs valeurs de $\lambda$. On choisi\textcolor{red}{t} par la suite la valeur de $\lambda$ qui minimise le taux d'erreur test.

\noindent Pour la r\'egression ridge, on a: 
\begin{verbatim}
> grille.lambda.ridge <- exp(seq(.5,14,0.1))
> reg.ridge.select <- cv.glmnet(x, y, alpha = 0, lambda=grille.lambda.ridge)
> log(reg.ridge.select$lambda.min)
[1] 2.2
\end{verbatim}

\noindent et donc, on choisi\textcolor{red}{t} $\lambda=e^{2.2}$.

\noindent Et pour la r\'egression lasso, on a :
\begin{verbatim}
> plot(reg.ridge.select)
> ## On selectionne un lambda en se basant sur la validation croisee
> 
> grille.lambda.lasso <- exp(seq(-3,6,0.1))
> 
> reg.lasso.select <- cv.glmnet(x,y,alpha=1,lambda=grille.lambda.lasso)
> log(reg.lasso.select$lambda.min)
[1] 3
\end{verbatim}

\noindent et donc, on choisi\textcolor{red}{t} $\lambda=e^{3}$.

\end{document}



